% !PRINTED 27
\weberpage{27}
%<<CONT
and from this, by (12) and (14), there is to be deduced:
\[
\begin{aligned}
\Phi' &= M\Phi + N\Psi,\\
\Psi' &= P\Phi + Q\Psi,
\end{aligned}
\]
in which the coefficients $M$, $N$, $P$, $Q$ are still to be determined.
We can carry this out by calculation. Without calculation the result
is obtained in the following way:

The $M$, $N$ are integral functions of the second degree in $p$, $q$ and
of the third degree in $m$, $n$. If one takes $p:q = m:n$, and so
$r = mq - np = 0$, then by (18) $\Phi'$ and $\dfrac{\partial \Phi'}{\partial p}$
pass over into linear combinations of $\varphi$ and $\psi$, and
therefore vanish. From this it follows that $M$ and $N$ are divisible
by $r^2$, and that the quotients are linear functions of $m$ and $n$. But
since for the two cases $\begin{pmatrix} m & n\\ p & q \end{pmatrix} = \begin{pmatrix} 1 & 0\\ 0 & 1 \end{pmatrix},\ \begin{pmatrix} 0 & 1\\ 1 & 0 \end{pmatrix}$
$\Phi'$ must pass over into $\Phi$ and into $\Psi$, and since the same
consideration can be applied to $\Psi'$, it follows:
\weq{19}{\begin{aligned}\Phi' &= r^2(m\Phi+n\Psi),\\ \Psi' &= r^2(p\Phi+q\Psi).\end{aligned}}
Thus (apart from the factor $r^2$) the functions
$\Phi$, $\Psi$ are transformed \gesperrt{cogrediently} with the $\varphi$, $\psi$.
(On the rigorous justification of these inferences see Bd. I, §\,20.)

From these results we can, in a simple way and always under the
assumption $\varphi = 0$, $\psi = 0$, express the $y_i^2$ by means of $\Phi$ and
$\Psi$, that is, solve the linear equations (12), (14).
For if we apply the formulae (19) to the substitution
\[
\begin{pmatrix} m, & n\\ p, & q \end{pmatrix} = \begin{pmatrix} \beta_1, & -\alpha_1\\ 0, & 1 \end{pmatrix}
\]
then $r = \beta_1$, the function $\psi'$ remains
unchanged $= \psi$, and $\varphi'$ arises from $\varphi$ when
\[
\alpha_1,\quad \alpha_2,\quad \alpha_3,\quad \alpha_4
\]
are replaced by
\[
0,\quad (\alpha_2\beta_1),\quad (\alpha_3\beta_1),\quad (\alpha_4\beta_1)
\]
Thus by (14)
\[
\Phi' = \beta_1^2(\alpha_2\beta_1)(\alpha_3\beta_1)(\alpha_4\beta_1)\,y_1^2
\]
and consequently by (19):
\weq{20}{\begin{aligned}
(\alpha_2\beta_1)(\alpha_3\beta_1)(\alpha_4\beta_1)\,y_1^2 &= \alpha_1\Phi - \beta_1\Psi,\\
(\alpha_1\beta_2)(\alpha_3\beta_2)(\alpha_4\beta_2)\,y_2^2 &= \alpha_2\Phi - \beta_2\Psi,\\
(\alpha_1\beta_3)(\alpha_2\beta_3)(\alpha_4\beta_3)\,y_3^2 &= \alpha_3\Phi - \beta_3\Psi,\\
(\alpha_1\beta_4)(\alpha_2\beta_4)(\alpha_3\beta_4)\,y_4^2 &= \alpha_4\Phi - \beta_4\Psi.
\end{aligned}}
